\clearpage
\fancyhead[L]{\small ZWEITER KOMET DES JAHRES 1805.}
\hypertarget{native-280-lower}{}
\begin{center}
ZWEITER KOMET DES JAHRES 1805.

\rule{0.6\linewidth}{0.4pt}

Monatliche Correspondenz. Band XIV. S. 181–186. August 1806.

\rule{0.6\linewidth}{0.4pt}
\end{center}

\begin{center}
\textsc{Gauss an von Zach.}
\end{center}

\begin{flushright}
Braunschweig, 8. Juli 1806.
\end{flushright}

\noindent --- --- --- Zum Kometen von 1805 habe ich seit meinem früheren Bericht an Sie nur noch die Beobachtungen von \textsc{Bouvard}, die Sie mir freundlicherweise mitgeteilt haben, mit meinen Elementen verglichen, sowohl mit den parabolischen (M.\,C., Juliheft 1806, \href{https://www.digitale-sammlungen.de/en/view/bsb10538606?page=87}{S.~77}) [\hyperlink{native-276}{S.~271} dieses Bandes] als auch mit den elliptischen (ebd., \href{https://www.digitale-sammlungen.de/en/view/bsb10538606?page=89}{S.~79}) [\hyperlink{native-277}{S.~272} dieses Bandes]. Der Vergleich sieht so aus: [Die Differenz ist R---B.]

\begin{center}\begingroup
\fontsize{8}{10}\selectfont
\setlength{\tabcolsep}{1pt}\renewcommand{\arraystretch}{1.12}
\begin{tabular}{l||c|c|c|c||c|c|c|c}
\shortstack{Zeit der Beob.} & \multicolumn{2}{c|}{\shortstack{In der Parabel\\berechnete}} & \multicolumn{2}{c||}{\shortstack{Differenz in}} & \multicolumn{2}{c|}{\shortstack{In der Ellipse\\berechnete}} & \multicolumn{2}{c}{\shortstack{Differenz in}} \\
Paris 1805. & Länge & Breite & Länge & Breite & Länge & Breite & Länge & Breite \\
\hline
Nov. 16.\,45203 & $29^\circ\,21'\,30''$ & $30^\circ\,2'\,32''$ & $-2'\,7''$ & $-8''$ & $29^\circ\,22'\,48''$ & $30^\circ\,1'\,18''$ & $-0'\,49''$ & $-1'\,22''$ \\
17.\,36386 & $28\,51\,5$ & $29\,49\,40$ & $-0\,18$ & $+48$ & $28\,51\,55$ & $29\,48\,47$ & $+0\,32$ & $-0\,4$ \\
18.\,39927 & $28\,14\,28$ & $29\,32\,42$ & $+0\,10$ & $-2$ & $28\,14\,43$ & $29\,32\,8$ & $+0\,25$ & $-0\,36$ \\
23.\,32241 & $24\,41\,33$ & $27\,25\,5$ & $+0\,29$ & $-30$ & $24\,39\,54$ & $27\,25\,31$ & $-1\,10$ & $-0\,4$ \\
30.\,51095 & $15\,43\,40$ & $19\,25\,32$ & $+4\,0$ & $+4$ & $15\,42\,3$ & $19\,26\,26$ & $+2\,23$ & $+0\,58$ \\
Dec. 5.\,29581 & $2\,7\,19$ & $3\,20\,5$ & $+0\,8$ & $-40$ & $2\,7\,20$ & $3\,21\,26$ & $+0\,9$ & $+0\,40$ \\
\end{tabular}\endgroup\end{center}

\par\indent Dieser Vergleich zeigt, dass die Ellipse mit diesen Beobachtungen nicht ganz so gut übereinstimmt wie mit den Beobachtungen aus Marseille. Zugleich scheint sich daraus zu ergeben, 1) dass die Beobachtungen aus Marseille besser sind und 2) dass auch bei den anderen die Parabel größere Differenzen ergibt als die Ellipse. Dabei muss ich noch bemerken, dass die hier für die Ellipse angesetzten Differenzen eigentlich nicht die wahren sind und noch erheblich geändert werden müssen. Die \emph{berechneten} Längen und Breiten sind hier nämlich frei von Parallaxe und Aberration. Ich habe sie mit den beobachteten verglichen, nachdem diese von Parallaxe und Aberration befreit worden waren (so wie Sie sie mir mitgeteilt haben). Diese bereits vorgenommene Reduktion hängt aber von den Abständen des Kometen von der Erde ab. Die Reduktion wegen der Parallaxe ist umgekehrt proportional zum Abstand, die wegen der Aberration dagegen direkt proportional. Die Abstände fallen in der Parabel und der Ellipse jedoch sehr verschieden aus. Zum Beispiel:

\begin{center}\begingroup
\fontsize{9}{11}\selectfont
\setlength{\tabcolsep}{2pt}\renewcommand{\arraystretch}{1.12}
\begin{tabular}{rlll}
Abstand in der Parabel: & Nov. 16 $=0.2489$, & Nov. 23 $=0.1761$, & Dec. 5 $=0.0638$ \\
in der Ellipse: & $=0.1794$ & $=0.1269$ & $=0.0460$ \\
\end{tabular}\endgroup\end{center}
\SourceLeaf{281}%
\par\indent Die Reduktion wegen der Parallaxe, die hier besonders beträchtlich ist, fällt für die Ellipse also ungefähr im Verhältnis 7:5 größer aus als für die Parabel, die Reduktion wegen der Aberration dagegen im selben Verhältnis kleiner. Da die beiden Reduktionen bei den Beobachtungen nicht einzeln angegeben sind und man aus den Angaben nur ihre Summe erschließen kann, und auch das nur bei den letzten vier Beobachtungen, müsste man die Reduktionen für die Ellipse neu berechnen. Diese Mühe habe ich mir nicht gemacht. Für die Parabel werden die Reduktionen dagegen so ausfallen, wie \textsc{Legendre} sie angenommen hat, da meine parabolischen Elemente so wenig von seinen abweichen. Schon durch bloßes Abschätzen erkennt man aber, dass die Beobachtung vom 5. Dezember, die einzige, die eine Meridianbeobachtung zu sein scheint, nach genauerer Reduktion noch viel besser mit den elliptischen Elementen übereinstimmen würde. Sollte ich \textsc{Bouvard}s ursprüngliche Rektaszensionen und Deklinationen noch zu Gesicht bekommen, werde ich sie auf dieselbe Weise mit meinen Elementen vergleichen wie die von \textsc{Thulis}. Vorerst habe ich diese Rechnungen beiseitegelegt, da sich über den Kometen nichts Bestimmteres sagen lassen wird, solange keine Beobachtungen bekannt werden, die den Zeitraum noch erweitern.

\par\indent \textsc{Legendre}s Werk, das Sie bei dieser Gelegenheit erwähnen, habe ich noch nicht gesehen. Ich habe mich absichtlich nicht darum bemüht, um bei der Arbeit an meiner Methode ganz in der Kette meiner eigenen Gedanken zu bleiben. Ein paar Worte, die \textsc{de la Lande} in der jüngsten \emph{Histoire de l'Astronomie}, 1805, fallen lässt, \emph{méthode des moindres quarrés}, bringen mich auf die Vermutung, dass ein Grundsatz, den ich schon seit zwölf Jahren bei mancherlei Rechnungen verwende und den ich auch in meinem Werk verwenden werde, obwohl er für meine Methode nicht wesentlich ist, --- dass dieser Grundsatz auch von \textsc{Legendre} verwendet worden ist. Im Übrigen dürfte es nur wenige Berührungspunkte geben, da unsere Arbeiten sehr verschiedene Gegenstände haben. Die Bestimmung einer Kometenbahn, bei der man eine Parabel voraussetzt, und die Bestimmung einer Planetenbahn, bei der man von jeder Hypothese absehen soll, erfordern eine ganz unterschiedliche Behandlung. Dem letzteren Problem ist mein Werk \emph{eigentlich} gewidmet, obwohl es auch zum ersteren vieles enthalten wird, was von mir selbst stammt.

\par\indent Ich bin übrigens überaus froh, dass ich die Methode, mit der ich die Bahnen von Ceres und Pallas berechnet hatte, nicht schon 1802 veröffentlicht habe, obwohl mich so viele dazu aufgefordert hatten. Denn seitdem habe ich weiter daran gearbeitet, die Methode selbst zu verbessern, besonders im vergangenen Winter, und in ihrer jetzigen Gestalt ähnelt sie ihrer ersten fast gar nicht mehr. Um davon einen Eindruck zu geben, will ich nur \emph{einen} Umstand erwähnen. Da das Problem so verwickelt ist, liegt es in der Natur der Sache, dass man bei der allerersten Annäherung einige Voraussetzungen machen muss, die nur näherungsweise richtig sind (wie zum Beispiel bei der Methode von Dr.~\textsc{Olbers} die Voraussetzung, dass die Sehnen bei der Erde und dem Kometen durch die mittleren Radiusvektoren im Verhältnis der Zeitabstände geschnitten werden). Voraussetzungen dieser Art liegen also auch bei meiner Methode der ersten Annäherung zugrunde. Sie weichen umso weniger von der Wahrheit ab, je kürzer die Zeitabstände sind. Man darf also die Beobachtungen, auf die man die erste Annäherung stützt, zeitlich nicht allzu weit auseinander wählen, weil man sonst wegen der nur näherungsweise richtigen Voraussetzung bei der ersten Rechnung zu weit von der Wahrheit entfernt bliebe und die Korrekturverfahren zu oft und mit zu großer Mühe wiederholen müsste. Doch konnte ich mit meiner Methode in ihrer Gestalt von 1802 bei Pallas gleich Beobachtungen verwenden, die 27 Tage auseinanderlagen. Viel weiter hätte ich allerdings nicht gehen wollen. Jetzt ist meine Methode dagegen so beschaffen, dass ich neulich, als ich die Beobachtungen von \textsc{Oriani} aus dem Jahr 1805, die Sie mir freundlicherweise mitgeteilt haben, für ein
\SourceLeaf{282}%
\emph{Beispiel} in meinem Werk verwenden wollte und sie also so behandeln musste, als wüsste ich von der Pallasbahn noch gar nichts, sogleich und mit dem besten Erfolg die äußersten, 71 Tage auseinanderliegenden Beobachtungen zugrunde legen konnte. Es besteht gar kein Zweifel, dass ich dabei noch beträchtlich weiter hätte gehen können.

